% !PRINTED 282
\weberpage{282}

\begin{center}
\S.\,72.

\gesperrt{Constitution of the Group $L_7$ of Degree 168.}
\end{center}

Among the simple groups found here, the next after the icosahedral group, which recurs here as $L_5$, is the group $L_7$ of degree 168, which, as we have seen, contains an octahedral group as a divisor. Since this remarkable group will meet us several times more later on, we wish here to enter somewhat more closely into its structure.

In \S.\,57 the octahedral group is represented by three elements $\chi, \omega, \Theta$ in the form:

\weq{1}{\chi^\lambda \omega^u \Theta^v, \quad \lambda = 0, 1, 2; \ \mu = 0, 1; \ \nu = 0, 1, 2, 3.} % [sic] formula prints Latin u, v; the parameter list prints Greek mu, nu

Between these elements there subsist the relations

\weq{2}{\omega\chi = \chi^2\omega, \quad \Theta\omega = \omega\Theta^3, \quad \Theta\chi = \chi^2\omega\Theta^2, \quad \Theta^2\chi = \chi\omega\Theta^3,}

and these conditions, when the degrees 3, 2, 4 of the elements $\chi, \omega, \Theta$ are added to them, we have shown to be sufficient to characterise the system (1) as an octahedral group.

From (2) we have derived in \S.\,57, as consequences, the formulae:

\weq{3}{\begin{gathered}
\omega\chi = \chi^2\omega, \quad \omega\chi^2 = \chi\omega, \quad \Theta^2\chi^2 = \chi^2\omega\Theta,\\
\Theta\chi = \chi^2\omega\Theta^2, \quad \Theta^2\chi = \chi\omega\Theta^3, \quad \Theta^3\chi = \chi^2\Theta,\\
\Theta\omega = \omega\Theta^3, \quad \Theta^2\omega = \omega\Theta^2, \quad \Theta^3\omega = \omega\Theta, \quad \Theta\chi^2 = \chi\Theta^3,
\end{gathered}}

which we shall make use of several times later.

In \S.\,71 we have also represented the octahedral group as a congruence group with respect to the modulus 7, and in the formulae (11) there we have found

\weq{4}{\chi = \begin{pmatrix} 0, & 2 \\ 3, & 1 \end{pmatrix}, \ \omega = \begin{pmatrix} 3, & -1 \\ 3, & -3 \end{pmatrix}, \ \Theta = \begin{pmatrix} 2, & 1 \\ 3, & 2 \end{pmatrix}}

For what follows, however, it is expedient, in the interest of a simpler calculation, to further trans-
%CONT>>

\bigskip
\hrule
\bigskip

%<<CONT
\gesperrt{Serret} (Cours d'alg\`ebre supérieure) and by C. \gesperrt{Jordan} (Traité des Substitutions). Cf. \gesperrt{Weber}, Elliptische Functionen etc., \S.\,84 f. The complete enumeration of all divisors is due to \gesperrt{Gierster} (Math. Ann., Bd. XVIII). A detailed treatment of the congruence groups in ``\gesperrt{Klein-Fricke}, Vorlesungen über die Theorie der elliptischen Modulfunctionen'', Bd. I, Leipzig 1890.
